\section{The partner field}
\label{ch:partner-field}

This chapter treats the second $2$-string field $\Upsilon_2$, the partner field, which makes the action local at order $\gs^4$ with elementary vertices only, by cancelling the exchange term that eliminating $\Sigma_2$ produces.
It expands the paper's appendix on the partner field, reproduced in full as the first block: the localization of inverse vertices, the partner at order $\gs^4$, the alternative where $\Sigma_3$ is introduced first, and the higher orders, with the count of couplings at order $\gs^5$, its deficit, and the number of partners through order $\gs^8$.
The second block is the older presentation from the pedagogical notes, which starts from the analogy with the induced action of two-dimensional gravity and discusses the split signature of the partner sector and its algebraic structure; this material is not in the paper.
The section between the two blocks, written for this document, gives the argument that limits the structures an elementary presentation can produce and shows that one partner suffices at order $\gs^5$, with the coupling $\vertex{5}(\Upsilon_2, \Sigma_3)$ absent.

\subsection{The partner field in the paper}
\label{sec:partner:paper}

The block states four results, with the following status.

\begin{itemize}
  \item The partner at order $\gs^4$: the action~\eqref{p:eq:wsp:action} with elementary vertices reproduces $S_{\le 6}$ through order $\gs^4$, the ratio of the coefficients $2$ and $1$ being forced by the matching, and $\Upsilon_2$ is real with a positive coefficient of $\vertex{4}(\Upsilon_2^2)$.
    \status{mixed}\record{017-chi-partner-liouville (R24)}
    The closure through order $\gs^4$ and the reality are proven at tree level, for generic off-shell $\Psi$, with the quartic form invertible on the quotient and the quintic and sextic vertices vanishing on its kernel, and, for the reality, the reality of the vertices; the closure was also checked numerically in a finite-dimensional model.
    The quantum status of the negative quadratic term of the partner is open.

  \item The count at order $\gs^5$: with elementary vertices and the fields $(\Psi, \Sigma_2, \Upsilon_2, \Sigma_3)$, three couplings meet three structures in a diagonal system, $\Sigma_3$ is forced, and one partner suffices.
    \status{mixed}\record{038-chi-scheme-ogs5-count (R42)}
    The count is a combinatorial argument, proven under the hypotheses below (\cref{sec:partner:strand-locality}), with the coupling $\vertex{5}(\Upsilon_2, \Sigma_3)$ absent; the closure coefficients of the table are numerical, from a random-tensor model on three draws, and one falsification leg of that computation, the branch with $\vertex{5}(\Upsilon_2, \Sigma_3)$ switched on, was left unsettled.

  \item The coefficients of the deficit~\eqref{p:eq:wsp:deficit}: $- \frac{42}{25}$ on $\mathcal I_1$ and $\frac{7 (2 + c_6)}{5}$ on $\mathcal I_2$, and $- 1$ on $\vertex{7}(\Psi^7)$.
    \status{Conditional(PH14 [$O(\gs^4)$, $O(\gs^5)$], PH15, PH16) -- Proven}\record{048-chi-scheme-ogs5-matching (R49)}
    The coefficients are derived from the combinatorics of the factors $1 / n!$ and measured in the model on three draws ($- 1.680$, $- 1.674$, $- 1.680$ for $- \frac{42}{25}$); the block quotes the numerical side only.

  \item The rule $\Sigma_2 \to \Sigma_2 + \Upsilon_2$, which identifies the partner with the multiplier $\Theta_2$ of the multiplier form (\cref{p:sec:all:multiplier}) and cancels every exchange term with a $\sprod{2}{2}^{-1}$ line at every order, at tree level and with the quartic form invertible on the quotient, is an exact identity, by bilinearity of $\vertex{4}$.
    \status{Proven}
    The number of partners through order $\gs^8$, $\Upsilon_3$ being needed from order $\gs^6$ if and only if $2 a_7 + b_7 \neq 0$, and the closure with $\Upsilon_2$ alone along the completed squares, were found in the finite-dimensional model.
    \status{Conditional(PH14, PH15, PH16) -- Verified in a finite model}\record{101-partner-count (no ledger entry; tags as proposed in the record)}
\end{itemize}

The hypotheses named in the tags are, in words: the exchange terms are not linear combinations of the elementary vertices and boundary terms, at the orders indicated; the negative quadratic term of the partner has no physical consequence, which is safe at tree level and open beyond; and the working hypothesis, which these results support, that elementary couplings of the partner cancel the exchange terms order by order.

The block is the paper's appendix as it stood before the paper was shortened, in the paper's notation: $\Upsilon_2$ for the partner, $\vertex{5 \cdot 5}$ for the exchange vertex, $\mathcal I_1$, $\mathcal I_2$ for the two exchange structures at order $\gs^5$.

\subsection{A partner field for the exchange term}
\label{p:app:wrong-sign}



This appendix gives the construction announced in \cref{p:sec:hs:sextic}: a second $2$-string auxiliary field $\Upsilon_2$ which makes the action at order $\gs^4$ local, with elementary vertices only, that is, the vertices $\vertex{n}$ of the original action with fields in their slots.
  Eliminating $\Sigma_2$ from the quintic-order action produces at order $\gs^4$ the exchange term $\vertex{5 \cdot 5}(\Psi^6)$ of~\eqref{p:eq:hs5:eliminate}, which the original action does not contain.
  For generic vertices, in which the exchange term is not a linear combination of the elementary vertices, the elementary vertices alone cannot cancel it, since it contains $\sprod{2}{2}^{-1}$.
  With the fields $(\Psi, \Sigma_2, \Sigma_3)$, each sextic coupling, $\vertex{6}(\Sigma_2^3)$, $\vertex{6}(\Sigma_3, \Sigma_2, \Psi)$ or $\vertex{6}(\Sigma_3^2)$, reduces on the classical values of the auxiliaries to a multiple of $\vertex{6}(\Psi^6)$ at this order.
    One structure is available to match two, the contact term $\vertex{6}(\Psi^6)$ and the exchange term, and no choice of coefficients closes the matching.
  The triangular form of \cref{p:sec:hs:sextic} places it inside the reduced vertex $\dvertex{6}$.


\subsubsection{Localizing inverse vertices}


In field theory, inverse operators arise in effective actions, and we know how to make them local: by the Gaussian identity, a field $X$ entering the action as $- K(X, X) + 2 J(X)$, with $K$ a quadratic form and $J$ a linear form, integrates out to $J K^{-1} J$.
  At tree level, integrating out $X$ amounts to replacing it by its classical value $X^* = K^{-1} J$.
  Read backwards, the identity writes a non-local term $J K^{-1} J$ as a local action with one more field.

An explicit example is two-dimensional gravity~\cite{Polyakov:1981:QuantumGeometryBosonic}.
  The Coulomb gas action and its value on the solution $\phi^* = - \square^{-1} R$ are:
  \begin{equation}
  \label{p:eq:wsp:coulomb}
  S
    = \int \dd^2 x \, \sqrt{g} \, \left(
      \frac{1}{2} \, \phi \square \phi
      + R \phi
    \right),
  \qquad
  S[\phi^*]
    = - \frac{1}{2} \int \dd^2 x \, \sqrt{g} \, R \square^{-1} R,
\end{equation}
  where $R[g]$ is the Ricci scalar of the metric $g$ and $\square$ its Laplacian.
  The second expression is the non-local induced gravitational action, and the first one makes it local by introducing $\phi$.


\subsubsection{Partner field at order \texorpdfstring{$\gs^4$}{gs4}}

We can apply the same procedure to decompose the composite vertex $\vertex{5 \cdot 5}$.
  At this point, we have two choices: either we decompose it after introducing $\Sigma_3$, or before.
  We start with the second option, which needs a single new field.

We add a $2$-string field $\Upsilon_2 \in \Hn{2} / \ker \sprod{2}{2}$, with $\Ngh(\Upsilon_2) = 4$ and the parity of $\Sigma_2$, since it fills two slots of the quartic vertex; it is a field in its own right, distinct from the partner field of ghost number five of \cref{p:sec:gauge}.
  We take $X = \Upsilon_2$ in the identity, with:
  \begin{equation}
  K(\Upsilon_2, \Upsilon_2)
    := - \frac{\gs^2}{4!} \, \vertex{4}(\Upsilon_2^2),
  \qquad
  J(\Upsilon_2)
    := \frac{\gs^3}{5!} \, \vertex{5}(\Upsilon_2, \Sigma_2, \Psi).
\end{equation}
  This gives the action:
  \begin{equation}
  \label{p:eq:wsp:partner-action}
  S_{\Upsilon}
    := \frac{\gs^2}{4!} \, \vertex{4}(\Upsilon_2^2)
      + \frac{2 \gs^3}{5!} \, \vertex{5}(\Upsilon_2, \Sigma_2, \Psi).
\end{equation}
  We write $\vertex{4}(\Upsilon_2, \Upsilon_2) = \pair[2]{\Upsilon_2}{\sprod{2}{2}(\Upsilon_2)}$ and $\vertex{5}(\Upsilon_2, \Sigma_2, \Psi) = \pair[2]{\Upsilon_2}{\sprod{3}{2}(\Sigma_2, \Psi)}$.
  The equation of motion of $\Upsilon_2$, normalized as the one of $\Sigma_2$ in \cref{p:sec:hs:quartic}, and its solution then read:
  \begin{equation}
  \label{p:eq:wsp:saddle}
  \mathcal{E}^{(4)}_{\Upsilon_2}
    := \frac{\gs}{5} \, \sprod{3}{2}(\Sigma_2, \Psi)
      + \sprod{2}{2}(\Upsilon_2)
    = 0,
  \qquad
  \Upsilon_2^{(4)*}
    = - \frac{\gs}{5} \, \sprod{2}{2}^{-1} \circ \sprod{3}{2}(\Sigma_2, \Psi).
\end{equation}
  There is no source $\vertex{4}(\Upsilon_2, \Psi^2)$, so the classical value of $\Upsilon_2$ starts at order $\gs$.
  Inserting it, we get $J K^{-1} J$:
  \begin{equation}
  \label{p:eq:wsp:eliminated}
  S_{\Upsilon}(\Upsilon_2^{(4)*})
    = - \frac{\gs^4}{600} \, \vertex{5}\big(\Sigma_2, \Psi, \sprod{2}{2}^{-1} \circ \sprod{3}{2}(\Sigma_2, \Psi)\big).
\end{equation}
  On the classical value $\Sigma_2 = \Psi^2 + O(\gs)$ of~\eqref{p:eq:hs5:saddle}, this is $- \frac{\gs^4}{600} \, \vertex{5 \cdot 5}(\Psi^6) + O(\gs^5)$, the exchange term with coefficient $- 1$.
  In the Feynman rules of \cref{p:sec:hs:quartic}, an $\Upsilon_2$ line carries the propagator $- \frac{12}{\gs^2} \, \sprod{2}{2}^{-1}$, opposite in sign to the one of a $\Sigma_2$ line.

The action through order $\gs^4$ is then:
  \begin{equation}
  \label{p:eq:wsp:action}
  \begin{aligned}
    \caction[\le 6]^{\Upsilon}
      &
      = \frac{1}{2} \, \pair{\Psi}{Q_B \Psi}
        + \frac{\gs}{3!} \, \vertex{3}(\Psi^3)
        - \frac{\gs^2}{4!} \, \vertex{4}(\Sigma_2^2)
        + \frac{2 \gs^2}{4!} \, \vertex{4}(\Sigma_2, \Psi^2)
        + \frac{\gs^3}{5!} \, \vertex{5}(\Sigma_2^2, \Psi)
      \\ &
      \qquad
        + \frac{\gs^2}{4!} \, \vertex{4}(\Upsilon_2^2)
        + \frac{2 \gs^3}{5!} \, \vertex{5}(\Upsilon_2, \Sigma_2, \Psi)
        + \frac{\gs^4}{6!} \, \vertex{6}(\Sigma_2^3).
  \end{aligned}
\end{equation}
  For the action explicit through order $\gs^4$, the reality condition of $\Upsilon_2$ has the sign of the one of $\Sigma_2$, whose strands it shares (\cref{p:app:reality}):
  \begin{equation}
  \Upsilon_2^\star
    = \Upsilon_2.
\end{equation}
  The action is cubic in the fields $(\Psi, \Sigma_2, \Upsilon_2)$, every vertex is elementary, and there is neither a reduced vertex nor a $\Sigma_3$ at this order.
  With $\Sigma_2$ and $\Upsilon_2$ at their classical values, it reproduces $S_{\le 6}$ through order $\gs^4$: the exchange term is cancelled by~\eqref{p:eq:wsp:eliminated}, and the sextic contact term is carried by $\vertex{6}(\Sigma_2^3)$.
  This holds at tree level, for generic off-shell $\Psi$, when $\sprod{2}{2}$ is invertible on the quotient and the quintic and sextic vertices vanish on its kernel.
  We checked it in a finite-dimensional model with generic symmetric tensors in place of the vertices (\cref{p:app:finite-models}).


\subsubsection{Partner of the kinetic term of \texorpdfstring{$\Sigma_3$}{Sigma3}}

In the first option, $\Sigma_3$ is introduced first, and the sextic action~\eqref{p:eq:hs6:S6tilde-explicit} contains two non-local terms:
  \begin{equation}
  \frac{\gs^4}{600} \, \vertex{5}\big(\Sigma_3, \sprod{2}{2}^{-1} \circ \sprod{3}{2}(\Sigma_3)\big)
    - \frac{2 \gs^4}{600} \, \vertex{5}\big(\Sigma_3, \sprod{2}{2}^{-1} \circ \sprod{3}{2}(\Psi, \Sigma_2)\big).
\end{equation}
  The first one, in the kinetic term of $\Sigma_3$, has the form $J K^{-1} J$ and is made local by a field $\Xi_2$:
  \begin{equation}
  \label{p:eq:wsp:local-kinetic}
  - \frac{\gs^4}{600} \, \vertex{4}(\Xi_2^2)
    + \frac{2 \gs^4}{600} \, \vertex{5}(\Xi_2, \Sigma_3),
  \qquad
  \Xi_2^*
    = \sprod{2}{2}^{-1} \circ \sprod{3}{2}(\Sigma_3).
\end{equation}
  The second one, in the coupling of $\Sigma_3$, cannot be made local by the same field: coupling $\Xi_2$ to $\Sigma_3 - \Psi \otimes \Sigma_2$ would produce, after elimination, the term $\vertex{5}\big(\Psi, \Sigma_2, \sprod{2}{2}^{-1} \circ \sprod{3}{2}(\Psi, \Sigma_2)\big)$, quartic in the fields, which the completed square has removed.
This option therefore needs more new fields than the second one, as do the intermediate-field representations of~\cite{Rivasseau:2010:LoopVertexExpansion,Lionni:2018:NoteIntermediateField}, which need about $3 n$ auxiliary fields to decompose an interaction of degree $2 n$.
  Decomposing the exchange term before introducing $\Sigma_3$ has the further advantage that the kinetic term of $\Sigma_3$ contains only $\vertex{6}$.
  Below, with a single partner field and a choice of couplings, the exchange terms with a $\sprod{2}{2}^{-1}$ line cancel exactly at every order, and the finite-dimensional model confirms that every term of the action, the kinetic terms of $\Sigma_3$ and $\Sigma_4$ included, remains an elementary vertex through order $\gs^8$.


\subsubsection{Higher orders}

Consider a coupling built on $\vertex{n}$ with fields of $s_1, \dots, s_m$ strands in its slots, where $s = 1$ for $\Psi$ and $s = k$ for a $k$-string field.
  By ghost-number saturation, it exists only for $\sum_i s_i = n$, and it carries $\gs^{n - 2}$.
Once the auxiliary fields are eliminated, it contributes first at the order:
  \begin{equation}
  \label{p:eq:wsp:order}
  \Big( \sum_i s_i - 2 \Big)
    + \sum_i o_i,
\end{equation}
  where $o_i$ is the power of $\gs$ at which the classical value of the field in slot $i$ starts: $o = 0$ for $\Psi$, $\Sigma_2$ and $\Sigma_3$, and $o = 1$ for $\Upsilon_2$ by~\eqref{p:eq:wsp:saddle}.
  At the order given by~\eqref{p:eq:wsp:order}, a coupling with $o = 0$ in every slot gives a contact term, since these fields reduce to powers of $\Psi$.
  A field with $o = 1$ brings one $\sprod{2}{2}^{-1}$ per slot and one extra power of $\gs$.
  On $(\Psi, \Sigma_2, \Upsilon_2)$, a cubic coupling with $o = 0$ in every slot fills at most six strands, so no coupling produces the contact term $\vertex{7}(\Psi^7)$ at order $\gs^5$.
  Hence $\Sigma_3$ is needed.
  The sextic contact term can then be shared between $\vertex{6}(\Sigma_2^3)$ and the quadratic and source terms of $\Sigma_3$, and the sextic terms of~\eqref{p:eq:wsp:action} become:
  \begin{equation}
  \label{p:eq:wsp:sextic-split}
  \frac{\gs^4}{6!} \, c_6 \, \vertex{6}(\Sigma_2^3)
    + \frac{\gs^4}{6!} \, (1 - c_6) \Big[
      - \vertex{6}(\Sigma_3^2)
      + 2 \, \vertex{6}(\Sigma_3, \Sigma_2, \Psi)
    \Big].
\end{equation}
  The split does not change the eliminated action at order $\gs^4$, but it enters at order $\gs^5$.

At order $\gs^5$, the eliminated action must contain $\vertex{7}(\Psi^7)$ and cancel the two exchange terms generated by the lower orders, $\mathcal{I}_1$ with three quintic vertices joined by two $\Sigma_2$ lines and $\mathcal{I}_2$ with a sextic and a quintic vertex joined by one:
  \begin{equation}
  \label{p:eq:wsp:S12}
  \mathcal{I}_1
    := \vertex{5}\big(\sprod{2}{2}^{-1} \circ \sprod{3}{2}(\Psi^3), \sprod{2}{2}^{-1} \circ \sprod{3}{2}(\Psi^3), \Psi\big),
  \qquad
  \mathcal{I}_2
    := \vertex{6}\big(\sprod{2}{2}^{-1} \circ \sprod{3}{2}(\Psi^3), \Psi^4\big),
\end{equation}
  On the classical values, the action with the sextic terms~\eqref{p:eq:wsp:sextic-split} differs from $S_{\le 7}$ at this order by a deficit:
  \begin{equation}
  \label{p:eq:wsp:deficit}
  7! \, \Big[ \caction^{\Upsilon}(\text{classical values}) - S_{\le 7} \Big]_{\gs^5}
    = - \vertex{7}(\Psi^7)
      - \frac{42}{25} \, \mathcal{I}_1
      + \frac{7 (2 + c_6)}{5} \, \mathcal{I}_2.
\end{equation}
  By~\eqref{p:eq:wsp:order}, three couplings are available at this order up to the choice of slots, and each produces exactly one of the three structures:
  \begin{center}
    \begin{tabular}{lll}
      \hline
      structure            & coupling                                                      & coefficient   \\
      \hline
      $\vertex{7}(\Psi^7)$ & $\frac{\gs^5}{7!} \, a_7 \, \vertex{7}(\Sigma_3^2, \Psi)$     & $a_7 = 1$     \\
      $\mathcal{I}_1$      & $\frac{\gs^3}{5!} \, c \, \vertex{5}(\Upsilon_2^2, \Psi)$     & $c = 1$       \\
      $\mathcal{I}_2$      & $\frac{\gs^4}{6!} \, d \, \vertex{6}(\Upsilon_2, \Sigma_2^2)$ & $d = 2 + c_6$ \\
      \hline
    \end{tabular}
  \end{center}
  The couplings $\vertex{7}(\Sigma_3, \Sigma_2^2)$ and $\vertex{6}(\Upsilon_2, \Sigma_3, \Psi)$ produce the same structures and can replace the first and the third one.
  The system is diagonal and none of the three couplings acts at order $\gs^4$, so the matching closes.
  The coupling $\vertex{5}(\Upsilon_2, \Sigma_3)$, allowed by the counting, must be absent: once $\Sigma_3$ has its kinetic term, it gives the quadratic form of $\Upsilon_2$ a term through the inverse of $\sprod{3}{3}$ at order $\gs^4$, which spoils the cancellation at that order.
  The deficit~\eqref{p:eq:wsp:deficit}, with no further structure in it, and the coefficients of the table were obtained in the finite-dimensional model, on three random draws.
  With these coefficients, the residual scales as $\gs^6$.
  One partner field thus suffices through order $\gs^5$, at the price of two couplings, one quadratic in $\Upsilon_2$ and one quadratic in $\Sigma_2$.
  The price is also structural: with $\vertex{6}(\Upsilon_2, \Sigma_2^2)$, eliminating $\Upsilon_2$ leaves an action quartic in $\Sigma_2$, and the auxiliary fields are no longer eliminated one at a time by exact Gaussian integrals.

The couplings of $\Upsilon_2$ in~\eqref{p:eq:wsp:action} and in the table follow a rule: replace $\Sigma_2$ by $\Sigma_2 + \Upsilon_2$ in every coupling with a $\Sigma_2$ leg, and keep the quadratic terms, with $+ \vertex{4}(\Upsilon_2^2)$ beside $- \vertex{4}(\Sigma_2^2)$ and no $\vertex{4}(\Upsilon_2, \Psi^2)$.
  It gives the coefficient $2$ of~\eqref{p:eq:wsp:partner-action} and $c = 1$, and at order $\gs^4$ the couplings $2 (1 - c_6) \, \vertex{6}(\Upsilon_2, \Sigma_3, \Psi)$, $3 c_6 \, \vertex{6}(\Upsilon_2, \Sigma_2^2)$, $3 c_6 \, \vertex{6}(\Upsilon_2^2, \Sigma_2)$ and $c_6 \, \vertex{6}(\Upsilon_2^3)$.
  By bilinearity of $\vertex{4}$, the action obtained in this way from an action $S[\Psi, \Sigma_2, \Sigma_3, \ldots]$ with elementary vertices and no partner is exactly:
  \begin{equation}
  \label{p:eq:wsp:rule}
  \caction^{\Upsilon}
    = S[\Psi, \Sigma_2 + \Upsilon_2, \Sigma_3, \ldots]
      + \frac{2 \gs^2}{4!} \, \vertex{4}(\Upsilon_2, \Sigma_2 + \Upsilon_2 - \Psi^2).
\end{equation}
  In the variables $\Sigma_2 + \Upsilon_2$ and $\Upsilon_2$, the partner field enters linearly, and the last term is the multiplier term of~\eqref{p:eq:all:multiplier-action} for $k = 2$, with $\Theta_2 = \frac{2 \gs^2}{4!} \, \sprod{2}{2}(\Upsilon_2)$.
  The partner field is thus the multiplier $\Theta_2$ of the all-orders form, written in the variables $(\Sigma_2, \Upsilon_2)$ of~\eqref{p:eq:wsp:partner-action}.
  Its elimination sets $\Sigma_2 + \Upsilon_2 = \Psi^2$ exactly, so every exchange term with a $\sprod{2}{2}^{-1}$ line cancels, at every order, for every $c_6$ and every choice of the couplings of $\Sigma_3, \Sigma_4, \ldots$.

The remaining action is the action of $\Sigma_3, \Sigma_4, \ldots$ at $\Sigma_2 = \Psi^2$.
  Its exchange terms contain no $\sprod{2}{2}^{-1}$ and are reached by no coupling of $\Upsilon_2$, and they decide how many partner fields are needed.
  At order $\gs^6$, the order of $\vertex{8}$, there is one such term, the exchange of $\Sigma_3$ between two couplings on $\vertex{7}$, with a coefficient proportional to $(2 a_7 + b_7)^2$ for the couplings $a_7 \, \vertex{7}(\Sigma_3^2, \Psi) + b_7 \, \vertex{7}(\Sigma_3, \Sigma_2^2)$.
  With the coupling of the table, $a_7 = 1$, a partner $\Upsilon_3$ of $\Sigma_3$ is needed from this order on, with the quadratic term $\vertex{6}(\Upsilon_3^2)$ and the source $2 \, \vertex{7}(\Upsilon_3, \Sigma_3, \Psi)$.
  If the field with the most strands is squared at the next orders as well, in $\vertex{8}(\Sigma_4, \Sigma_3, \Psi)$ and $\vertex{9}(\Sigma_4^2, \Psi)$, an $\Upsilon_4$ follows at order $\gs^8$, and the partners form a tower which is the multiplier form itself, one strand at a time.
  With couplings whose terms linear in $\Sigma_3 - \Psi^3$ and $\Sigma_4 - \Psi^4$ at the classical point do not contribute before order $\gs^8$, the remaining exchange terms vanish through that order and no other partner is needed.
  These couplings are $- \vertex{7}(\Sigma_3^2, \Psi) + 2 \, \vertex{7}(\Sigma_3, \Sigma_2^2)$ at order $\gs^5$, a completed square on $\Sigma_2 = \Psi^2$; $- \vertex{8}(\Sigma_4^2) + 2 \, \vertex{8}(\Sigma_4, \Sigma_2^2)$ at order $\gs^6$, which brings in $\Sigma_4$ at that order, as in \cref{p:sec:hs:octic}; $\lambda \, [\vertex{9}(\Sigma_4^2, \Psi) - 2 \, \vertex{9}(\Sigma_4, \Sigma_3, \Sigma_2) - \vertex{9}(\Sigma_3^3)] + \vertex{9}(\Sigma_3^3)$ at order $\gs^7$, whose term linear in $\Sigma_3 - \Psi^3$, proportional to $3 + \lambda$, contributes first at order $\gs^{10}$, so that $\lambda$ is free through order $\gs^8$; and $\rho \, \vertex{10}(\Sigma_4^2, \Sigma_2) + (1 - \rho) \, \vertex{10}(\Sigma_4, \Sigma_3^2)$ at order $\gs^8$.
  Through order $\gs^8$, the fields are then $\Psi$, $\Sigma_2$, $\Sigma_3$, $\Sigma_4$ and the single partner $\Upsilon_2$.
  The finite-dimensional model confirms this through order $\gs^8$ at $c_6 = 0$, and at the values of the rule also for $c_6 = \frac{1}{3}$: with $\Upsilon_2$ alone and these couplings the matching closes, and the only couplings of $\Upsilon_2$ that close it are those of the rule, up to splits between couplings producing the same exchange terms; the model detects the one it reaches, the departure of $d$ from the rule discussed below, at order $\gs^8$.
  Carrying the results of this subsection over to the string vertices requires, besides tree level and the invertibility of $\sprod{2}{2}$, $\sprod{3}{3}$ and $\sprod{4}{4}$ on their quotients, that the exchange terms are not linear combinations of the elementary vertices.

For $c_6 \neq 1$, the coupling $d \, \vertex{6}(\Upsilon_2, \Sigma_2^2)$ of the table, $d = 2 + c_6$, differs from the value $3 c_6$ of the rule.
  With the couplings above at $c_6 = 0$, it still closes the orders $\gs^5$ to $\gs^7$, since at order $\gs^5$ only the sum of the coefficients of $\vertex{6}(\Upsilon_2, \Sigma_2^2)$ and $\vertex{6}(\Upsilon_2, \Sigma_3, \Psi)$ enters, but it leaves at order $\gs^8$ a residue proportional to $d^2$ and $d^3$ on two exchange terms that no cubic coupling reaches.
  With the couplings of the table and their $\Upsilon_3$, order $\gs^8$ was not computed.
  Since such an $\Upsilon_3$ sets $\Sigma_3 = \Sigma_2 \otimes \Psi$ and the term in $d^3$, three $\Sigma_2$ lines on a single $\vertex{7}$, contains no $\Sigma_3$ line, we expect the same residue there.


\subsection{Strand locality and one partner at order \texorpdfstring{$\gs^5$}{gs5}}
\label{sec:partner:strand-locality}

The count of the block compares three couplings with three structures.
Two arguments lie behind it and are not in the paper: why no other structure can appear in the deficit, and why a partner of $\Sigma_3$ brings nothing at this order.
Both are statements about which contractions the elimination of the auxiliary fields can produce, at the level of algebra, for generic vertices.
They hold at tree level, with the quartic form and the quadratic form of $\Sigma_3$ invertible on their quotients, for exchange terms that are not linear combinations of the elementary vertices, and with the coupling $\vertex{5}(\Upsilon_2, \Sigma_3)$ set to zero, as the block requires.

\paragraph{Strand locality.}
In a presentation with elementary vertices only, every term of the action on the classical values is a tree contraction with whole slots.
Its nodes are elementary vertices $\vertex{n}$, its internal lines are the inverses $\sprod{2}{2}^{-1}$ on $\Hn{2}$ and $\sprod{3}{3}^{-1}$ on $\Hn{3}$, and its external lines are single strands of $\Psi$.
The reason is that each vertex of the action takes whole fields in its slots, the only single-strand field being $\Psi$, and the Gaussian elimination of an auxiliary field replaces a pair of its slots by the inverse of its quadratic form, an operator on the whole multi-string space, which joins one complete slot to one complete slot.
No step of the elimination routes the strands of one slot into two different vertices.

The consequence is that a strand-split contraction is never generated.
The exchange vertex $\vertex{5 \cdot 5}$, read as a trilinear form on $2$-string states, sends the two strands of its middle argument into the two different quintic vertices (\cref{p:eq:hs5:exchange-vertex}).
Its contraction at order $\gs^5$,
\begin{equation}
  \mathcal I_3
    := \vertex{5 \cdot 5}\big( \Psi^2, \Psi^2, \sprod{2}{2}^{-1} \circ \sprod{3}{2}(\Psi^3) \big),
\end{equation}
is therefore outside the reach of any elementary presentation.
It appears in the remainder at order $\gs^5$ of the minimal presentation, through the composite vertex $\dvertex{6} = \vertex{6} - \frac{6}{5} \vertex{5 \cdot 5}$, not through any elementary structure, and the partner presentation is never required to produce it.
An enumeration of the whole-slot trees of degree seven with the fields $(\Psi, \Sigma_2, \Upsilon_2, \Sigma_3)$ then gives exactly the three structures of the block, $\vertex{7}(\Psi^7)$, $\mathcal I_1$ and $\mathcal I_2$, and no fourth one, in particular none with $\sprod{3}{3}^{-1}$.
The model confirms the absence of $\mathcal I_3$ and of any fourth structure in the deficit.

This statement concerns the contractions a Feynman expansion can produce; it says nothing about whether a kernel can be realized on the worldsheet.

\paragraph{The coupling $\vertex{5}(\Upsilon_2, \Sigma_3)$.}
The enumeration needs this coupling to vanish.
It fills five strands with one partner leg, so by~\eqref{p:eq:wsp:order} it is admissible at order $\gs^4$.
At order $\gs^5$ it would mix $\Upsilon_2$ and $\Sigma_3$ at a power of $\gs$ that is marginal between their quadratic terms, so that chains alternating $\sprod{2}{2}^{-1}$ and $\sprod{3}{3}^{-1}$ lines through further insertions are all of the same order and resum into the inverse of the joint quadratic form; the linear count of the block would not apply to that branch.
Setting it to zero costs no freedom at the lower orders, and the block gives the reason it must be absent: it spoils the cancellation at order $\gs^4$ once $\Sigma_3$ has its quadratic term.
What the count would become with this coupling switched on was not settled.

\paragraph{One partner suffices.}
Consider a partner $\Upsilon_3$ of $\Sigma_3$, with a quadratic term at order $\gs^4$.
Suppose first it has no source of order $\gs^4$, so that its classical value starts at order $\gs$.
By~\eqref{p:eq:wsp:order}, a coupling containing $\Upsilon_3$ that contributes at order $\gs^5$ or below fills at most $5 - 1 + 2 = 6$ strands; it is then linear in $\Upsilon_3$ with at most three further strands, which makes it a source of $\Upsilon_3$ at order $\gs^4$ or below, for example $\vertex{6}(\Upsilon_3, \Sigma_2, \Psi)$.
This contradicts the assumption, so every coupling of such a partner contributes first at order $\gs^6$, and $\Upsilon_3$ is irrelevant at order $\gs^5$.
If instead $\Upsilon_3$ is given the source $\vertex{6}(\Upsilon_3, \Sigma_2, \Psi)$, its classical value starts at order one and it is a second $\Sigma_3$ with the opposite sign of the quadratic term: it only re-splits the sextic contact term and changes the coefficients of the deficit, without a new direction.
The one remaining coupling, the mixing $\vertex{5}(\Upsilon_3, \Upsilon_2)$, sources $\Upsilon_3$ at order $\gs^4$ with a composite classical value, and its exchange contributes a structure with $\sprod{3}{3}^{-1}$ at order $\gs^4$, which is absent from the matching at that order; it must vanish, excluded by the lower order and not by order $\gs^5$.
Hence the partners do not have to mirror the tower of the $\Sigma_k$ at order $\gs^5$; the block shows how this continues, with $\Upsilon_3$ needed from order $\gs^6$ on one route and not on the route of completed squares.

\subsection{The Liouville analogy in the pedagogical notes}
\label{sec:partner:liouville}

The second block is the section of the pedagogical notes in which the partner was found.
It writes the partner as $\chi$ and works with the same action as the paper's~\eqref{p:eq:wsp:action}, in the notation of the notes; the box below lists the differences that occur in it.
It starts from the exact rewriting of the quintic-order action as a Gaussian in the deviation of $\Sigma$ from its classical value, which the notes give in their section on gauge transformations:
\begin{equation}
  \label{eq:partner:S3rewrite}
  S_3
    = S^{(5)}_{\mathrm{eff}}
      - \frac{\gs^2}{4!} \, K_\Psi(\tilde X, \tilde X)
      + \frac{\gs^4}{600} \, v K_\Psi^{-1} v,
  \qquad
  K_\Psi
    := V_4 - \frac{\gs}{5} \, V_5(\cdot, \cdot, \Psi),
\end{equation}
with $v := V_5(\cdot, \Psi^{\otimes 2}, \Psi)$ and the deviation $\tilde X := \Sigma - \Psi^{\otimes 2} - \frac{\gs}{5} \, K_\Psi^{-1} v$, so that the equation of $\Sigma$ is $\tilde X = 0$.
Here $S_3$ is the cubic action through the quintic order and $S^{(5)}_{\mathrm{eff}}$ the original action through the same order.

\begin{quote}
  \small
  \textbf{Notation of the block.}
  \begin{itemize}
    \item $\chi$ is the partner field $\Upsilon_2$; $\Sig$ is $\Sigma_2$, and $\Sym^2 \Hilb_0$ is the $2$-string space $\Hn{2}$ of the paper, on the quotient by the kernel of the quartic form.
    \item $V_n$ is $\vertex{n}$; the kinetic term $V_4(\Sigma, \Sigma)$ is $\vertex{4}(\Sigma_2^2) = \pair[2]{\Sigma_2}{\sprod{2}{2}(\Sigma_2)}$, and $V_4^{-1}$ is $\sprod{2}{2}^{-1}$; $\mtwo$ is $\sprod{2}{2}$.
    \item $\Psi^{\otimes 2}$, also written $P$, is $\Psi^2$; the source $v = V_5(\cdot, \Psi^{\otimes 2}, \Psi)$ is $\sprod{3}{2}(\Psi^3)$ read through the $2$-string pairing, and the feedback $v V_4^{-1} v$ is $\vertex{5 \cdot 5}(\Psi^6)$.
    \item The composite $\mathcal C$ is the exchange vertex $\vertex{5 \cdot 5}$, and $W_6 = V_6 - \frac{6}{5} \mathcal C$ is the reduced vertex $\dvertex{6}$; $\tilde M$ is the dressed quadratic form of $\Sigma_3$.
    \item $(a_\chi, b_\chi) = (1, 2)$ are the coefficients of $\vertex{4}(\Upsilon_2^2)$ and $\vertex{5}(\Upsilon_2, \Sigma_2, \Psi)$ in~\eqref{p:eq:wsp:partner-action}; $S_{\mathrm{HS}}$ is the cubic action at quartic order.
    \item Scheme C is the partner form of the all-orders discussion (\cref{ped:sec:allorders:chi}).
    \item Reality: the notes write $\Sigma^\ddagger = \Sigma$, the paper $\Sigma_k^\star = (-1)^k \, \Sigma_k$ and $\Psi^\star = - \Psi$. For odd $k$ the two conventions differ by a sign; for the $2$-string fields of this block, $\Sigma_2$ and $\Upsilon_2$, they agree, and $\Upsilon_2^\star = \Upsilon_2$.
  \end{itemize}
\end{quote}

The block leaves three questions open at its end: the quantum status of the negative direction, the gauge law and exact rewriting of the partner sector, and the extension to order $\gs^5$.
The last one is answered in the paper's appendix above, and by \cref{sec:partner:strand-locality}: one partner suffices through order $\gs^5$, with the coupling $\vertex{5}(\Upsilon_2, \Sigma_3)$ absent, and, in a finite model with the couplings of the completed squares, through order $\gs^8$.
The first two remain open.
The light-cone form of the split-signature sector, the twist lemma and the coloured coalgebra of the partner are not in the paper.

\subsubsection{The Liouville analogy and the wrong-sign partner
\texorpdfstring{$\chi$}{chi}}
\label{ped:sec:milestones:chi}

The feedback $vK_\Psi^{-1}v$ has a familiar shape: it is the induced nonlocal action of
2d gravity, $S\sim R\,\Delta^{-1}R$, obtained by integrating the Liouville field out of
$S=\phi\Delta\phi+R\phi$. The analogy is exact, and
running it \emph{backwards} produces a new presentation of the sextic
(checked numerically in a finite-dimensional model, see below).

\subparagraph{The dictionary (exact).} In the exact rewriting \eqref{eq:partner:S3rewrite} the
deviation $\tilde X$ is a Gaussian field with kernel $K_\Psi$ and linear source
$\propto v$; integrating it out \emph{produces} the nonlocal composite:
\begin{center}
	\begin{tabular}{lcl}
		\toprule
		2d gravity & & here \\
		\midrule
		Liouville field $\phi$ & $\leftrightarrow$ & deviation $\tilde X=\Sig-\Sig_*(\Psi)$ \\
		kinetic operator $\Delta$ & $\leftrightarrow$ & dressed kernel $K_\Psi$ \\
		curvature source $R$ (metric-dep.) & $\leftrightarrow$ & source $v=V_5(\cdot,\Psi^{\otimes2},\Psi)$ ($\Psi$-dep.) \\
		induced $R\,\Delta^{-1}R$ & $\leftrightarrow$ & feedback $\tfrac{\gs^4}{600}\,vK_\Psi^{-1}v$ \\
		\bottomrule
	\end{tabular}
\end{center}

\subparagraph{The constructive converse.} In 2d gravity one also \emph{localises} the
nonlocal Polyakov action by re-introducing $\phi$. Here the object to remove is the
feedback \emph{overshoot}: it must be \emph{subtracted} at $O(\gs^4)$ --- which is
exactly why the minimal presentation of the sextic order needs the composite $-\tfrac65\mathcal C$ inside $W_6$, and
the presentation with $\Sig_3$ the composite $\tilde M$. A subtraction is the induced action of a
\emph{wrong-sign} Gaussian field. Introduce a second bilocal auxiliary
$\chi\in\Sym^2\Hilb_0$ ($\Ngh(\chi)=4$, even --- the same quantum numbers as $\Sig$)
with \emph{elementary vertices only}:
\begin{equation}
	S_\chi
	= +\,\frac{\gs^2}{4!}\,a_\chi V_4(\chi,\chi)
	+ \frac{\gs^3}{5!}\,b_\chi V_5(\chi,\Sig,\Psi),
	\label{ped:eq:Schi}
\end{equation}
--- note the \emph{plus} sign of the kinetic term, opposite to
$-\tfrac{\gs^2}{4!}V_4(\Sig,\Sig)$, and that every term is cubic in fields
($\chi\Sig\Psi$ is trilinear; the direct source $v(\chi)=V_5(\chi,\Psi^{\otimes2},\Psi)$
would be quartic in fields and is \emph{not} used). $\chi$ has no
$V_4(\chi,\Psi^{\otimes2})$ source, so it talks to the theory only through the
trilinear $O(\gs^3)$ coupling. On the $\Sig$-saddle ($\Sig\to P$, so
$V_5(\chi,\Sig,\Psi)\to v(\chi)$) the $\chi$-sector is a wrong-sign sourced Gaussian
whose on-shell value is
\begin{equation}
	S_\chi\big|_{\chi_*}
	= -\,\frac{b_\chi^2}{4a_\chi}\cdot\frac{(\gs^3/5!)^2\,4!}{\gs^2}\;vV_4^{-1}v+O(\gs^5)
	= -\,\frac{b_\chi^2}{a_\chi}\cdot\frac{\gs^4}{2400}\,vV_4^{-1}v+O(\gs^5),
	\label{ped:eq:chivalue}
\end{equation}
so cancelling the feedback $+\tfrac{\gs^4}{600}vV_4^{-1}v$ \emph{forces}
\begin{equation}
	\boxed{\;\frac{b_\chi^2}{a_\chi}=4,\qquad (a_\chi,b_\chi)=(1,2)\;}
	\label{ped:eq:chiforced}
\end{equation}
--- the universal HS $1{:}2$ kinetic-to-source ratio, reappearing once more. The result
is an all-elementary-vertex cubic action that closes at $O(\gs^4)$ with \emph{no}
composite, \emph{no} dressed vertex, and \emph{no} $\Sig_3$:
\begin{equation}
	\boxed{\;
	S^{\chi}_4
	= S_{\rm HS}
	+ \frac{\gs^3}{5!}V_5(\Sig,\Sig,\Psi)
	+ \frac{\gs^2}{4!}V_4(\chi,\chi)
	+ \frac{2\gs^3}{5!}V_5(\chi,\Sig,\Psi)
	+ \frac{\gs^4}{6!}V_6(\Sig,\Sig,\Sig).
	\;}
	\label{ped:eq:S4chi}
\end{equation}
Machine check (generic non-$\Loo$ tensors, joint $(\Sig,\chi)$ saddle, Richardson in
$\gs$): without $\chi$ the mismatch/$\gs^4$ reproduces $+vV_4^{-1}v/600$ to
$3\times10^{-5}$; with $(a_\chi,b_\chi)=(1,2)$ it vanishes at the $10^{-3}$ level with
residual scaling exponent $\approx0.93$ (i.e.\ the mismatch is $O(\gs^5)$); $b_\chi=1$
leaves exactly $(1-b_\chi^2/4)$ of the residual; the 5-point matching is uncontaminated
($\chi$-exchange starts at $O(\gs^4)$).

\subparagraph{Consistency with the no-go for elementary vertices, and what $\chi$ is.} The no-go for elementary vertices at order $\gs^4$
had content $\Sig_2\oplus\Sig_3$ and its strong form (\enquote{no elementary
presentation whatsoever}) was explicitly a conjecture; the $\chi$-scheme settles that
strong form in the \emph{negative} --- \emph{provided indefinite-signature
(negative-kinetic) auxiliaries are admitted as an elementary presentation}: the
falsification inherits the tree-level-only unitarity caveat below as part of its
scope. The mechanism is instructive: a $\Sig_3$-exchange is
$V_6$-built and can never reach the feedback (this is the no-go), while the
$\chi$-exchange is $V_5V_4^{-1}V_5$-built --- the \emph{same} composite as the feedback
--- with a sign-flipped propagator. The impossible trade $V_4^{-1}\to V_6^{-1}$ is
replaced by the possible pairing $+V_4^{-1}\to\pm V_4^{-1}$. As for what $\chi$
\emph{is}: its quantum numbers ($\Ngh=4$, even) show it is \emph{not} the odd
up-component $\Sig_\uparrow$ of an unconstrained two-string field (that sits at
$\Ngh=3$) and \emph{not} the ghost-5 partner $\Pi$ of the gauge analysis at quartic order --- at the parity
level it is genuinely a new field. But the \emph{structure} it creates is telling: the
joint $(\tilde X,\tilde\chi)$ quadratic form is split-signature,
$-K_\Psi\oplus+V_4$ --- equal to $-K_\Psi\oplus+K_\Psi$ only at leading order in
the dressing, since no $V_5(\chi,\chi,\Psi)$ coupling exists and the $\chi$-kinetic
stays bare --- i.e.\ at leading order purely off-diagonal in the light-cone
combinations $\Pi_\pm=\tilde X\mp\tilde\chi$:
\begin{equation}
	-K_\Psi(\tilde X,\tilde X)+K_\Psi(\tilde\chi,\tilde\chi)
	=-K_\Psi(\Pi_+,\Pi_-),
	\label{ped:eq:lightcone}
\end{equation}
the shape of a BV quartet / contractible-pair sector with a $V_4$-type (non-$\hat Q$)
off-diagonal kernel --- suggestive of an already-gauge-fixed enlarged theory, i.e.\ a
$V_5$-sourced cousin of the open question whether the quartic Hubbard--Stratonovich action is the gauge fixing of a larger theory. \emph{Caveats, honest:}
$\chi$ carries negative kinetic energy. At tree level this is harmless in the precise
sense that the kinetic operator is a contact kernel, the
telescoping of the $\Sig$ exchange to a contact term applies verbatim with a sign, no propagating
negative-norm states arise, and the formal Gaussian is signature-indifferent
--- but the quantum status of the negative direction, the
$\chi$-sector gauge law and exact rewriting, and the $O(\gs^5)$ extension (one $\chi$,
or a wrong-sign tower mirroring the $\Sig_k$?) are open.

\subparagraph{The algebraic home: sign-twisted $\IBL$ and the coloured coalgebra.} Does
the $\chi$-action still carry an $\IBL$ structure? Two-part answer: the $\chi$-sector
alone is a \emph{sign-twisted copy} of the same $\IBL$ data --- this is a theorem, by a
two-line counting argument --- while the \emph{coupled} $(\Psi,\Sig,\chi)$ system lives
on the symmetric coalgebra of \emph{several copies} of $\Hilb_0$, a coloured structure
whose systematic development is the natural next step.

\emph{(a) The twist lemma.} Every structure map is built from one geometric bracket by
$\Delta$-raisings and $\Delta$-self-gluings, and each costs exactly one $\Delta$:
starting from $\ell_n$ ($n$ inputs, $1$ output) and applying $r$ raisings ($-1$ input,
$+1$ output) and $s$ self-gluings ($-2$ inputs, $+1$ genus), one lands on the signature
$(k,\ell,g)=(n-r-2s,\,1+r,\,s)$ with $\Delta$-count
\begin{equation}
	\#\Delta\big(\pklg\big)=r+s=\ell-1+g.
	\label{ped:eq:deltacount}
\end{equation}
This count is \emph{conserved} by the coderivation composition: gluing $j$ legs of a
$(k_2,\ell_2,g_2)$-operation into a $(k_1,\ell_1,g_1)$-operation produces the signature
$(k_1{+}k_2{-}j,\,\ell_1{+}\ell_2{-}j,\,g_1{+}g_2{+}j{-}1)$, whose count is
\begin{equation}
	(\ell_1{+}\ell_2{-}j)-1+(g_1{+}g_2{+}j{-}1)
	=(\ell_1{-}1{+}g_1)+(\ell_2{-}1{+}g_2)
	\label{ped:eq:deltaconserved}
\end{equation}
--- the $j$-dependence cancels. Hence \emph{every} component relation of
$\hat{\mathfrak p}\circ\hat{\mathfrak p}=0$ at fixed signature is homogeneous in the
total $\Delta$-count, and the rescaling
\begin{equation}
	\Delta\to\lambda\Delta,
	\qquad
	\pklg\to\lambda^{\ell-1+g}\,\pklg
	\qquad(\lambda\neq0)
	\label{ped:eq:lambdatwist}
\end{equation}
maps $\IBL$ data to $\IBL$ data at all signatures and genera. (For $\lambda$ a formal
parameter this is the familiar $\hbar$-type grading freedom; the new use is)
$\lambda=-1$: the wrong-sign sector. The $\chi$-kinetic $+V_4(\chi,\chi)$ is
$G_2(\chi,\mtwo^-\chi)$ with $\mtwo^-=(-\Delta)\circ\ell_3=-\mtwo$, i.e.\ the
$\chi$-copy uses the $\lambda=-1$ twist of the \emph{same} products built from the
\emph{same} geometric vertices --- it satisfies the same $\IBL$ relations and imposes no
new condition. This is the algebraic reason the wrong-sign trick is free.

\emph{(b) The coloured coalgebra.} The coupled system no longer fits in
$\IBL(S\Hilb_0)$: the symmetric coalgebra of one copy of $\Hilb_0$ has exactly
\emph{one} two-string slot, and $(\Sig,\chi)$ are two independent $\Sym^2$-valued
fields. The minimal home is an $\IBL$ structure on several
copies of $\Sym\Hilb$: the symmetric coalgebra of a \emph{coloured} complex,
\begin{equation}
	E\big(\Hilb_0^{(1)}\oplus\Hilb_0^{(2)}\big)\;\cong\;E\big(\Hilb_0^{(1)}\big)\otimes
	E\big(\Hilb_0^{(2)}\big),
	\label{ped:eq:colouredE}
\end{equation}
with $\Sig$ the $\Sym^2$-component of the first factor and $\chi$ of the second, the
two copies identified as vector spaces but carrying opposite twists
($\lambda=+1$ vs $-1$, equivalently pairings $\pm\omega_1^{\otimes2}$), and the master
field \emph{truncated}: no single-string colour-2 component, no mixed-strand
components. The geometric vertices are colour-blind ($V_5(\chi,\Sig,\Psi)$ is one
trinion with legs of three colours); the colour enters \emph{only} through the
propagator signs on colour-2 edges --- the coloured-operad statement of the scheme with colour-2 propagators
(\cref{ped:sec:allorders:chi}). Two concrete facts support the truncation's consistency at
this order: the $\chi$-saddle is exactly minus the feedback shift of the $\Sig$-saddle,
\begin{equation}
	\Sig_*=P+\tfrac{\gs}5 K_\Psi^{-1}v+O(\gs^2),
	\qquad
	\chi_*=-\tfrac{\gs}5 V_4^{-1}v+O(\gs^2)
	\quad\Longrightarrow\quad
	\Sig_*+\chi_*=P+O(\gs^2),
	\label{ped:eq:saddlesum}
\end{equation}
so the coloured master field returns to the Veronese variety at leading order; and the
quadratic sector is invariant under a rigid $O(1,1)$ boost $\delta\tilde X=
\varepsilon\tilde\chi$, $\delta\tilde\chi=\varepsilon\tilde X$ of the split-signature
pair \eqref{ped:eq:lightcone} --- the first structure in the whole HS tower admitting
\emph{any} independent two-string transformation, albeit rigid and broken already
at $O(\gs^3)$ by the $V_5$ couplings ($\delta V_5(\chi,\Sig,\Psi)=\varepsilon
[V_5(\Sig,\Sig,\Psi)+V_5(\chi,\chi,\Psi)]\neq0$; the $K_\Psi$-vs-$V_4$ kinetic
asymmetry breaks it at $O(\gs)$ within the quadratic sector itself); whether a field-dependent
completion survives is part of the same open question on a larger theory. \emph{Gauge transformations of all
fields}: on \eqref{ped:eq:colouredE} the universal coderivation law
$\delta\Phi=\hat{\mathfrak p}(\Lambda,\e^\Phi)$ applies verbatim, predicting
saddle-tracking laws for \emph{both} $\Sig$ and $\chi$
($\delta\chi=\delta[\chi_*(\Psi)]+\text{homogeneous rotation}$, the
contractible-pair class with $\chi_*$ from \eqref{ped:eq:saddlesum}); the consistency
condition is that this law \emph{preserves the truncation} (does not source the
single-string colour-2 or mixed components). The twist lemma beyond
$\ell\le2$ on the toy model of the $\Delta$-raised products, the truncation-preservation, and the $\delta\chi$
law have not been verified.


